%%=============================================================================
%% Voorwoord
%%=============================================================================

\chapter*{\IfLanguageName{dutch}{Woord vooraf}{Preface}}%
\label{ch:voorwoord}

%% TODO:
%% Het voorwoord is het enige deel van de bachelorproef waar je vanuit je
%% eigen standpunt (``ik-vorm'') mag schrijven. Je kan hier bv. motiveren
%% waarom jij het onderwerp wil bespreken.
%% Vergeet ook niet te bedanken wie je geholpen/gesteund/... heeft

Één van mijn grootste passies is voetbal, daarom ging ik bij het zoeken van een onderwerp op zoek naar een opportuniteit binnen deze sector. Ik dacht echter dat de kans dat ik iets ging vinden relatief klein was. 

Toen ik bij KAA Gent de kans kreeg om dit onderwerp te bespreken was ik meteen geïnteresseerd, niet enkel was het een interessant onderwerp in de sportsector maar ook binnen de club waar ik al jarenlang fan ben. Na een positieve bespreking met de club heb ik met mijn promotor de mogelijkheden besproken en zijn we samen tot een concrete onderzoeksvraag gekomen.

Graag wil ik dan ook mijn promotor, Tom Antjon, en co-promotor, Adriaan Martiny, bedanken voor deze opportuniteit. Tijdens de samenwerking werd er altijd duidelijk aangegeven wat goed liep en wat verbeterd kon worden, om zo het project in de juiste richting te sturen.

Ik zou graag ook mijn zus, vader en moeder bedanken voor het altijd steunen van mijn studies en bachelorproef, die altijd vroegen of ze ergens konden helpen en ook klaarstonden indien nodig.

Ik hoop dat dit onderzoek bijdraagt aan waardevolle inzichten in de mogelijkheden van dataverwerking in de sportwereld.